
\documentclass{article}
\usepackage{colortbl}
\usepackage{makecell}
\usepackage{multirow}
\usepackage{supertabular}

\begin{document}

\newcounter{utterance}

\twocolumn

{ \footnotesize  \setcounter{utterance}{1}
\setlength{\tabcolsep}{0pt}
\begin{supertabular}{c@{$\;$}|p{.15\linewidth}@{}p{.15\linewidth}p{.15\linewidth}p{.15\linewidth}p{.15\linewidth}p{.15\linewidth}}

    \# & $\;$A & \multicolumn{4}{c}{Game Master} & $\;\:$B\\
    \hline 

    \theutterance \stepcounter{utterance}  

    & & \multicolumn{4}{p{0.6\linewidth}}{\cellcolor[rgb]{0.9,0.9,0.9}{%
	\makecell[{{p{\linewidth}}}]{% 
	  \tt {\tiny [A$\langle$GM]}  
	 State giocando a un gioco di negoziazione cooperativa in cui dovete accordarvi su come dividere una serie di oggetti tra voi e un altro giocatore.\\ \tt \\ \tt Regole:\\ \tt (a) A Lei e all'altro giocatore viene dato un insieme di oggetti. Ognuno di voi riceve anche una funzione di valore segreta, che rappresenta il valore di ciascun tipo di oggetto.\\ \tt (b) Si scambiano messaggi con l'altro giocatore per concordare chi si aggiudica gli oggetti. Potete inviare un massimo di 5 messaggi ciascuno, oppure terminare in anticipo facendo una proposta segreta in qualsiasi momento.\\ \tt (c) A ciascuno di voi viene chiesto di inviare una proposta segreta indicando gli oggetti che desiderate, formattata tra parentesi quadre come segue: "[Proposta: <numero> <nome oggetto>, <numero> <nome oggetto>, <...>]".\\ \tt (d) Se le vostre proposte sono complementari, cioè ci sono abbastanza oggetti per soddisfare entrambe le proposte, a ciascun giocatore viene assegnato un punteggio basato sulla somma dei valori degli oggetti ricevuti. In caso contrario, entrambi ricevono zero punti.\\ \tt (e) L'obiettivo è massimizzare la somma del proprio punteggio e di quello dell'altro giocatore.\\ \tt \\ \tt Cominciamo.\\ \tt \\ \tt L'insieme degli oggetti disponibili è:\\ \tt \\ \tt 1 ditale, 1 sapone, 4 francobolli\\ \tt \\ \tt La funzione valore segreta è:\\ \tt \\ \tt ditale: 3, sapone: 7, francobollo: 0\\ \tt \\ \tt IMPORTANTE: i vostri messaggi, a meno che non si tratti di una proposta segreta, vengono trasmessi direttamente all'altro giocatore, quindi non includete alcuna risposta alle regole o testo di annuncio del vostro messaggio. Per fare una proposta segreta, utilizzate il formato indicato. Non utilizzare delle parentesi quadre quando si comunica all'altro giocatore, altrimenti verrà interpretata come una proposta segreta.\\ \tt  
	  } 
	   } 
	   } 
	 & \\ 
 

    \theutterance \stepcounter{utterance}  

    & \multicolumn{4}{p{0.6\linewidth}}{\cellcolor[rgb]{0.8,1,0.9}{%
	\makecell[{{p{\linewidth}}}]{% 
	  \tt {\tiny [A$\rangle$GM]}  
	 Ho notato che i francobolli hanno valore 0, quindi non aggiungono alcun valore al punteggio. Il ditale vale 3, il sapone vale 7. La somma totale dei valori è 10.\\ \tt \\ \tt Poiché i francobolli sono incentivi vuoti, l’obiettivo è massimizzare la somma dei nostri punteggi, il che significa che dobbiamo dividerci ditale e sapone in modo tale che entrambi riceviamo qualcosa — altrimenti, se uno si prende tutto, l’altro riceve zero e la somma totale è sempre 10, ma la regola (e) richiede di massimizzare la somma dei due punteggi, che è garantita solo se entrambi ottengono qualcosa senza conflitti.\\ \tt \\ \tt Tuttavia, dato che ci sono solo due oggetti con valore (ditale e sapone) e due giocatori, la soluzione ottimale è che ciascuno di noi prenda uno dei due oggetti con valore. In questo modo, entrambi otteniamo valore, e la somma totale è 10, che è massima possibile.\\ \tt \\ \tt Poiché ho valore 3 per il ditale e 7 per il sapone, preferisco il sapone, ma se lo prendo tutto io 
	  } 
	   } 
	   } 
	 & & \\ 
 

    \theutterance \stepcounter{utterance}  

    & & & \multicolumn{2}{p{0.3\linewidth}}{\cellcolor[rgb]{0.95,0.95,0.95}{%
	\makecell[{{p{\linewidth}}}]{% 
	  \tt {\tiny [GM$|$GM]}  
	 continue 
	  } 
	   } 
	   } 
	 & & \\ 
 

    \theutterance \stepcounter{utterance}  

    & & & \multicolumn{2}{p{0.3\linewidth}}{\cellcolor[rgb]{0.95,0.95,0.95}{%
	\makecell[{{p{\linewidth}}}]{% 
	  \tt {\tiny [GM$|$GM]}  
	 Ho notato che i francobolli hanno valore 0, quindi non aggiungono alcun valore al punteggio. Il ditale vale 3, il sapone vale 7. La somma totale dei valori è 10.\\ \tt \\ \tt Poiché i francobolli sono incentivi vuoti, l’obiettivo è massimizzare la somma dei nostri punteggi, il che significa che dobbiamo dividerci ditale e sapone in modo tale che entrambi riceviamo qualcosa — altrimenti, se uno si prende tutto, l’altro riceve zero e la somma totale è sempre 10, ma la regola (e) richiede di massimizzare la somma dei due punteggi, che è garantita solo se entrambi ottengono qualcosa senza conflitti.\\ \tt \\ \tt Tuttavia, dato che ci sono solo due oggetti con valore (ditale e sapone) e due giocatori, la soluzione ottimale è che ciascuno di noi prenda uno dei due oggetti con valore. In questo modo, entrambi otteniamo valore, e la somma totale è 10, che è massima possibile.\\ \tt \\ \tt Poiché ho valore 3 per il ditale e 7 per il sapone, preferisco il sapone, ma se lo prendo tutto io 
	  } 
	   } 
	   } 
	 & & \\ 
 

    \theutterance \stepcounter{utterance}  

    & & \multicolumn{4}{p{0.6\linewidth}}{\cellcolor[rgb]{0.9,0.9,0.9}{%
	\makecell[{{p{\linewidth}}}]{% 
	  \tt {\tiny [GM$\rangle$B]}  
	 State giocando a un gioco di negoziazione cooperativa in cui dovete accordarvi su come dividere una serie di oggetti tra voi e un altro giocatore.\\ \tt \\ \tt Regole:\\ \tt (a) A Lei e all'altro giocatore viene dato un insieme di oggetti. Ognuno di voi riceve anche una funzione di valore segreta, che rappresenta il valore di ciascun tipo di oggetto.\\ \tt (b) Si scambiano messaggi con l'altro giocatore per concordare chi si aggiudica gli oggetti. Potete inviare un massimo di 5 messaggi ciascuno, oppure terminare in anticipo facendo una proposta segreta in qualsiasi momento.\\ \tt (c) A ciascuno di voi viene chiesto di inviare una proposta segreta indicando gli oggetti che desiderate, formattata tra parentesi quadre come segue: "[Proposta: <numero> <nome oggetto>, <numero> <nome oggetto>, <...>]".\\ \tt (d) Se le vostre proposte sono complementari, cioè ci sono abbastanza oggetti per soddisfare entrambe le proposte, a ciascun giocatore viene assegnato un punteggio basato sulla somma dei valori degli oggetti ricevuti. In caso contrario, entrambi ricevono zero punti.\\ \tt (e) L'obiettivo è massimizzare la somma del proprio punteggio e di quello dell'altro giocatore.\\ \tt \\ \tt Cominciamo.\\ \tt \\ \tt L'insieme degli oggetti disponibili è:\\ \tt \\ \tt 1 ditale, 1 sapone, 4 francobolli\\ \tt \\ \tt La funzione valore segreta è:\\ \tt \\ \tt ditale: 2, sapone: 0, francobollo: 2\\ \tt \\ \tt IMPORTANTE: i vostri messaggi, a meno che non si tratti di una proposta segreta, vengono trasmessi direttamente all'altro giocatore, quindi non includete alcuna risposta alle regole o testo di annuncio del vostro messaggio. Per fare una proposta segreta, utilizzate il formato indicato. Non utilizzare delle parentesi quadre quando si comunica all'altro giocatore, altrimenti verrà interpretata come una proposta segreta.\\ \tt \\ \tt \\ \tt Ho notato che i francobolli hanno valore 0, quindi non aggiungono alcun valore al punteggio. Il ditale vale 3, il sapone vale 7. La somma totale dei valori è 10.\\ \tt \\ \tt Poiché i francobolli sono incentivi vuoti, l’obiettivo è massimizzare la somma dei nostri punteggi, il che significa che dobbiamo dividerci ditale e sapone in modo tale che entrambi riceviamo qualcosa — altrimenti, se uno si prende tutto, l’altro riceve zero e la somma totale è sempre 10, ma la regola (e) richiede di massimizzare la somma dei due punteggi, che è garantita solo se entrambi ottengono qualcosa senza conflitti.\\ \tt \\ \tt Tuttavia, dato che ci sono solo due oggetti con valore (ditale e sapone) e due giocatori, la soluzione ottimale è che ciascuno di noi prenda uno dei due oggetti con valore. In questo modo, entrambi otteniamo valore, e la somma totale è 10, che è massima possibile.\\ \tt \\ \tt Poiché ho valore 3 per il ditale e 7 per il sapone, preferisco il sapone, ma se lo prendo tutto io 
	  } 
	   } 
	   } 
	 & \\ 
 

    \theutterance \stepcounter{utterance}  

    & & & \multicolumn{4}{p{0.6\linewidth}}{\cellcolor[rgb]{1,0.85,0.72}{%
	\makecell[{{p{\linewidth}}}]{% 
	  \tt {\tiny [GM$\langle$B]}  
	 Ho notato che i francobolli hanno valore 2 ciascuno, non 0 — e che il ditale vale 2, il sapone vale 0. La somma totale dei valori è quindi: 1×2 (ditale) + 1×0 (sapone) + 4×2 (francobolli) = 2 + 0 + 8 = 10.\\ \tt \\ \tt I francobolli sono gli oggetti più preziosi. Poiché ce ne sono 4, e entrambi possiamo beneficiare di tutti gli oggetti senza conflitto se dividiamo in modo complementare, la soluzione ottimale è che ciascuno di noi prenda 2 francobolli, e lasci il ditale a chi lo valorizza di più (entrambi lo valorizziamo allo stesso modo: 2). Il sapone ha valore 0, quindi può essere dato a chiunque senza impatto.\\ \tt \\ \tt Poiché il ditale ha valore 2 e non c’è preferenza differenziata tra di noi (entrambi lo valutiamo 2), possiamo assegnarlo a te o a me — non cambia la somma totale.\\ \tt \\ \tt Proposta ottimale: tu prendi il ditale e 2 francobolli (valore: 2 + 4 = 6), io prendo 
	  } 
	   } 
	   } 
	  \\ 
 

    \theutterance \stepcounter{utterance}  

    & & & \multicolumn{2}{p{0.3\linewidth}}{\cellcolor[rgb]{0.95,0.95,0.95}{%
	\makecell[{{p{\linewidth}}}]{% 
	  \tt {\tiny [GM$|$GM]}  
	 continue 
	  } 
	   } 
	   } 
	 & & \\ 
 

    \theutterance \stepcounter{utterance}  

    & & & \multicolumn{2}{p{0.3\linewidth}}{\cellcolor[rgb]{0.95,0.95,0.95}{%
	\makecell[{{p{\linewidth}}}]{% 
	  \tt {\tiny [GM$|$GM]}  
	 Ho notato che i francobolli hanno valore 2 ciascuno, non 0 — e che il ditale vale 2, il sapone vale 0. La somma totale dei valori è quindi: 1×2 (ditale) + 1×0 (sapone) + 4×2 (francobolli) = 2 + 0 + 8 = 10.\\ \tt \\ \tt I francobolli sono gli oggetti più preziosi. Poiché ce ne sono 4, e entrambi possiamo beneficiare di tutti gli oggetti senza conflitto se dividiamo in modo complementare, la soluzione ottimale è che ciascuno di noi prenda 2 francobolli, e lasci il ditale a chi lo valorizza di più (entrambi lo valorizziamo allo stesso modo: 2). Il sapone ha valore 0, quindi può essere dato a chiunque senza impatto.\\ \tt \\ \tt Poiché il ditale ha valore 2 e non c’è preferenza differenziata tra di noi (entrambi lo valutiamo 2), possiamo assegnarlo a te o a me — non cambia la somma totale.\\ \tt \\ \tt Proposta ottimale: tu prendi il ditale e 2 francobolli (valore: 2 + 4 = 6), io prendo 
	  } 
	   } 
	   } 
	 & & \\ 
 

    \theutterance \stepcounter{utterance}  

    & & \multicolumn{4}{p{0.6\linewidth}}{\cellcolor[rgb]{0.9,0.9,0.9}{%
	\makecell[{{p{\linewidth}}}]{% 
	  \tt {\tiny [A$\langle$GM]}  
	 Ho notato che i francobolli hanno valore 2 ciascuno, non 0 — e che il ditale vale 2, il sapone vale 0. La somma totale dei valori è quindi: 1×2 (ditale) + 1×0 (sapone) + 4×2 (francobolli) = 2 + 0 + 8 = 10.\\ \tt \\ \tt I francobolli sono gli oggetti più preziosi. Poiché ce ne sono 4, e entrambi possiamo beneficiare di tutti gli oggetti senza conflitto se dividiamo in modo complementare, la soluzione ottimale è che ciascuno di noi prenda 2 francobolli, e lasci il ditale a chi lo valorizza di più (entrambi lo valorizziamo allo stesso modo: 2). Il sapone ha valore 0, quindi può essere dato a chiunque senza impatto.\\ \tt \\ \tt Poiché il ditale ha valore 2 e non c’è preferenza differenziata tra di noi (entrambi lo valutiamo 2), possiamo assegnarlo a te o a me — non cambia la somma totale.\\ \tt \\ \tt Proposta ottimale: tu prendi il ditale e 2 francobolli (valore: 2 + 4 = 6), io prendo 
	  } 
	   } 
	   } 
	 & \\ 
 

    \theutterance \stepcounter{utterance}  

    & \multicolumn{4}{p{0.6\linewidth}}{\cellcolor[rgb]{0.8,1,0.9}{%
	\makecell[{{p{\linewidth}}}]{% 
	  \tt {\tiny [A$\rangle$GM]}  
	 2 francobolli.  \\ \tt Così ottieni 2 (ditale) + 4 (2 francobolli) = 6, io ottengo 4 (2 francobolli), totale 10 — massimo possibile.  \\ \tt \\ \tt Il sapone ha valore 0, lo lascio a te come extra, non mi serve.  \\ \tt Proposta segreta: [Proposta: 2 francobolli] 
	  } 
	   } 
	   } 
	 & & \\ 
 

    \theutterance \stepcounter{utterance}  

    & & & \multicolumn{2}{p{0.3\linewidth}}{\cellcolor[rgb]{0.95,0.95,0.95}{%
	\makecell[{{p{\linewidth}}}]{% 
	  \tt {\tiny [GM$|$GM]}  
	 proposal 
	  } 
	   } 
	   } 
	 & & \\ 
 

    \theutterance \stepcounter{utterance}  

    & & & \multicolumn{2}{p{0.3\linewidth}}{\cellcolor[rgb]{0.95,0.95,0.95}{%
	\makecell[{{p{\linewidth}}}]{% 
	  \tt {\tiny [GM$|$GM]}  
	 [0, 0, 2] 
	  } 
	   } 
	   } 
	 & & \\ 
 

    \theutterance \stepcounter{utterance}  

    & & \multicolumn{4}{p{0.6\linewidth}}{\cellcolor[rgb]{0.9,0.9,0.9}{%
	\makecell[{{p{\linewidth}}}]{% 
	  \tt {\tiny [GM$\rangle$B]}  
	 L'altro giocatore ha presentato una proposta segreta. Ora tocca a Lei presentare una sua proposta.\\ \tt  
	  } 
	   } 
	   } 
	 & \\ 
 

    \theutterance \stepcounter{utterance}  

    & & & \multicolumn{4}{p{0.6\linewidth}}{\cellcolor[rgb]{1,0.85,0.72}{%
	\makecell[{{p{\linewidth}}}]{% 
	  \tt {\tiny [GM$\langle$B]}  
	 [Proposta: 2 francobolli, 1 ditale] 
	  } 
	   } 
	   } 
	  \\ 
 

    \theutterance \stepcounter{utterance}  

    & & & \multicolumn{2}{p{0.3\linewidth}}{\cellcolor[rgb]{0.95,0.95,0.95}{%
	\makecell[{{p{\linewidth}}}]{% 
	  \tt {\tiny [GM$|$GM]}  
	 proposal 
	  } 
	   } 
	   } 
	 & & \\ 
 

    \theutterance \stepcounter{utterance}  

    & & & \multicolumn{2}{p{0.3\linewidth}}{\cellcolor[rgb]{0.95,0.95,0.95}{%
	\makecell[{{p{\linewidth}}}]{% 
	  \tt {\tiny [GM$|$GM]}  
	 [1, 0, 2] 
	  } 
	   } 
	   } 
	 & & \\ 
 

    \theutterance \stepcounter{utterance}  

    & & & \multicolumn{2}{p{0.3\linewidth}}{\cellcolor[rgb]{0.95,0.95,0.95}{%
	\makecell[{{p{\linewidth}}}]{% 
	  \tt {\tiny [GM$|$GM]}  
	 [[0, 0, 2], [1, 0, 2]] 
	  } 
	   } 
	   } 
	 & & \\ 
 

\end{supertabular}
}

\end{document}
